% Created 2026-06-08 Mon 10:13
% Intended LaTeX compiler: pdflatex
\documentclass[11pt]{article}
\usepackage[utf8]{inputenc}
\usepackage[T1]{fontenc}
\usepackage{graphicx}
\usepackage{longtable}
\usepackage{wrapfig}
\usepackage{rotating}
\usepackage[normalem]{ulem}
\usepackage{amsmath}
\usepackage{amssymb}
\usepackage{capt-of}
\usepackage{hyperref}
\author{Elif Güler}
\date{\today}
\title{}
\hypersetup{
 pdfauthor={Elif Güler},
 pdftitle={},
 pdfkeywords={},
 pdfsubject={},
 pdfcreator={Emacs 30.2 (Org mode 9.7.11)}, 
 pdflang={English}}
\begin{document}

\tableofcontents

\section{MAPRO Full Analysis Test Run}
\label{sec:orgac3fa4d}

Command executed:

\begin{verbatim}
bash ../scripts/driver.sh \
    -n ../data/neidb \
    < ../data/test_list.txt
\end{verbatim}

The driver script processes each organism listed in \texttt{test\_list.txt} independently.
\subsection{1. Aquifex aeolicus VF5 (aae)}
\label{sec:orgec11320}

Output:

\begin{verbatim}
Analyzing aae Aquifex aeolicus VF5
Collecting 1 genome record
Downloading: tdata.zip
...
Error: There are no genome assemblies that match your query.
\end{verbatim}

Explanation:

\begin{itemize}
\item The type strain was found in the Neighbors database.
\item Only one target genome accession was available.
\item NCBI Datasets attempted to download the genome.
\item No complete/chromosome-level assembly was available.
\item As a result, no genome FASTA file was downloaded.
\end{itemize}

Consequences:

\begin{verbatim}
ls: cannot access 'aae/targets/'
ls: cannot access 'aae/neighbors'
\end{verbatim}

Since no genomes were available, MAPRO could not build a phylogeny and therefore skipped marker discovery for this species.

Result:

\begin{itemize}
\item Analysis failed.
\item No markers produced.
\item No primers produced.
\end{itemize}

---
\subsection{2. Agrobacterium fabrum C58 (afa)}
\label{sec:orga8a6c0a}

Output:

\begin{verbatim}
Collecting 20 genome records
Collecting 68 genome records
\end{verbatim}

Explanation:

\begin{itemize}
\item 20 genomes were identified as target genomes.
\item 68 genomes were identified as neighboring genomes.
\end{itemize}

The genomes were downloaded and rehydrated successfully.

Output:

\begin{verbatim}
Found 20 of 20 files for rehydration
Completed 20 of 20

Found 68 of 68 files for rehydration
Completed 68 of 68
\end{verbatim}

A phylogeny was constructed and a representative target genome was selected:

\begin{verbatim}
using tGCF_003667945.1_ASM366794v1_genomic.fna as target representative
\end{verbatim}

Marker discovery (Fur):

\begin{verbatim}
Step             Sequences   Length      Ns
Subtraction_1        5305    1487141      0
Intersection         2188     469635  18938
Subtraction_2          55      39423  26398
\end{verbatim}

Interpretation:

\begin{itemize}
\item Subtraction\textsubscript{1}:
sequences present in target genomes but absent from neighbors.

\item Intersection:
sequences shared among all target genomes.

\item Subtraction\textsubscript{2}:
final candidate marker regions.
\end{itemize}

Final result:

\begin{itemize}
\item 55 marker candidates.
\item Total marker length = 39,423 bp.
\end{itemize}

---
\subsection{3. Aliivibrio fischeri ES114 (afi)}
\label{sec:org45a1de2}

Output:

\begin{verbatim}
Collecting 12 genome records
Collecting 15 genome records
\end{verbatim}

Meaning:

\begin{itemize}
\item 12 target genomes.
\item 15 neighboring genomes.
\end{itemize}

Representative genome:

\begin{verbatim}
using tGCF_000241785.1_ASM24178v1_genomic.fna as target representative
\end{verbatim}

Marker discovery:

\begin{verbatim}
Subtraction_1       4484   3626854     1202
Intersection        6041   2001815   169339
Subtraction_2        363     98908    28173
\end{verbatim}

Interpretation:

\begin{itemize}
\item Large amount of target-specific sequence was found.
\item After filtering and intersection steps, 363 candidate markers remained.
\end{itemize}

Final result:

\begin{itemize}
\item 363 marker candidates.
\item Total marker length = 98,908 bp.
\end{itemize}

---
\subsection{4. Aeromonas hydrophila ATCC 7966 (ahy)}
\label{sec:orgbcd4cae}

Output:

\begin{verbatim}
Collecting 4 genome records
Collecting 5 genome records
\end{verbatim}

Meaning:

\begin{itemize}
\item 4 target genomes.
\item 5 neighboring genomes.
\end{itemize}

Representative genome:

\begin{verbatim}
using tGCF_000014805.1_ASM1480v1_genomic.fna as target representative
\end{verbatim}

Marker discovery:

\begin{verbatim}
Subtraction_1       1597   530113     0
Intersection         330    59775  4225
Subtraction_2         20    11874   776
\end{verbatim}

Interpretation:

\begin{itemize}
\item Much fewer marker regions survived filtering.
\item Only 20 candidate markers remained.
\end{itemize}

Final result:

\begin{itemize}
\item 20 marker candidates.
\item Total marker length = 11,874 bp.
\end{itemize}

---
\subsection{5. Amycolatopsis mediterranei U32 (ame)}
\label{sec:orgbd71a7f}

Output:

\begin{verbatim}
Collecting 5 genome records
Collecting 49 genome records
\end{verbatim}

Meaning:

\begin{itemize}
\item 5 target genomes.
\item 49 neighboring genomes.
\end{itemize}

Representative genome:

\begin{verbatim}
using tGCF_000196835.1_ASM19683v1_genomic.fna as target representative
\end{verbatim}

Marker discovery:

\begin{verbatim}
Subtraction_1      12457   4250144       0
Intersection       11520   4038681     200
Subtraction_2        539   1228442  643033
\end{verbatim}

Interpretation:

\begin{itemize}
\item Very large amount of sequence survived the first filtering step.
\item More than 1.2 Mb of candidate marker sequence remained.
\item This species appears to have many regions distinguishable from its neighbors.
\end{itemize}

Final result:

\begin{itemize}
\item 539 marker candidates.
\item Total marker length = 1,228,442 bp.
\end{itemize}

---
\subsection{General Workflow}
\label{sec:orgc181876}

For each organism, MAPRO performs:

\begin{enumerate}
\item Identify target and neighboring taxa using the Neighbors database.
\item Download genomes from NCBI.
\item Build a phylogeny.
\item Select a representative target genome.
\item Create a Fur database.
\item Run Fur marker discovery.
\item Produce marker sequences.
\item Design PCR primers from the discovered markers.
\end{enumerate}

The first organism (Aquifex aeolicus) failed because no suitable assembly was available from NCBI.

All subsequent organisms completed successfully and entered the marker-discovery stage.
\end{document}
